\documentclass[11pt]{article}

\usepackage[margin=1in]{geometry}
\usepackage{amsmath,amssymb,bm}
\usepackage{booktabs}
\usepackage{enumitem}
\usepackage{xurl}
\usepackage[hidelinks]{hyperref}

\setlist[itemize]{leftmargin=1.4em}
\setlist[enumerate]{leftmargin=1.6em}
\newcommand{\dd}{\mathrm{d}}
\newcommand{\jj}{\mathrm{j}}
\newcommand{\cc}{\mathrm{c.c.}}
\newcommand{\Heav}{\mathrm{H}}

\title{Piezoelectric Energy Harvester: Modeling Assumptions, ROM Derivation, and Slow-Flow Analysis}
\author{Working note for collaborators}
\date{\today}

\begin{document}
\maketitle

\section*{Purpose}

This note organizes the reasoning path behind the piezoelectric energy
harvesting model used in the weekly meeting slides and MATLAB simulations.  It
is not intended to be a polished paper derivation.  The goal is to make the
modeling choices, reduction steps, and analysis pipeline clear enough that a
collaborator can check the logic and connect the equations to the simulation
files.

\paragraph{Main sources.}
The derivation below is organized from the archived Word notes
\path{Old_sources/piezoelectric beam harvester.docx},
\path{Old_sources/ROM - send to Mohammod.docx},
\path{Old_sources/ROM - multiple - send to Mohammod.docx},
\path{Old_sources/Damping analysis.docx}, and the weekly meeting slides in
\path{docs/weekly_meeting/}.  The corresponding cleaned MATLAB archive is in
\path{simulations/}.

\section{Physical Picture and Assumptions}

\subsection{Beam and Electrical Paths}

The physical system is a cantilever piezoelectric beam under prescribed base
motion.  Two piezoelectric electrical paths are used:

\begin{enumerate}
  \item \textbf{Path 1: harvester path.}  This path is connected to a shunted
  electrical circuit.  It is the path from which harvested electrical energy is
  measured.  The harvester circuit contains capacitance, resistance, linear
  inductance, and an effective cubic nonlinear inductive term.
  \item \textbf{Path 2: modulator path.}  This path is not solved as an
  autonomous circuit in the reduced model.  Its voltage is prescribed and used
  as a controllable input.  In the main simulations it is taken near the base
  excitation frequency so that it slowly modulates the effective forcing.
\end{enumerate}

The first path is therefore a dynamical electrical degree of freedom.  The
second path is an externally imposed actuation/modulation signal.

\subsection{Mechanical Assumptions}

\begin{itemize}
  \item The beam is modeled as an Euler-Bernoulli beam.
  \item The total transverse displacement is decomposed into base motion plus
  relative beam deflection:
  \begin{equation}
    w(x,t)=w_b(t)+w_{\mathrm{rel}}(x,t).
  \end{equation}
  \item The base motion is prescribed.  The archived notes often write
  $w_b(t)=g(t)$, and for harmonic studies use
  \begin{equation}
    w_b(t)=w_0 e^{\jj \omega_b t}+\cc.
  \end{equation}
  \item In the initial ROM derivation, structural damping of the beam is
  neglected.  A later extension includes a modal damping ratio $\xi$.
  \item The first ROM uses one mechanical mode:
  \begin{equation}
    w_{\mathrm{rel}}(x,t)=\sum_r \eta_r(t)\phi_r(x)
    \approx \eta_1(t)\phi_1(x).
  \end{equation}
  \item The mode shape is mass-normalized:
  \begin{equation}
    \int_0^l m\phi_1^2(x)\,\dd x = 1.
  \end{equation}
\end{itemize}

\paragraph{Important correction.}
The earliest ROM used the analytic first bending mode of a plain cantilever.
The January 17 meeting identified that this approximation shifted the relevant
electromechanical frequency enough to cause disagreement with the full beam
simulation.  The corrected ROM uses a mode computed from the FEM eigenanalysis
with the piezoelectric/electrical coupling included.  In the cleaned MATLAB
archive this is the role of:

\begin{itemize}
  \item \texttt{simulations/fem/eigenanalysis\_corrected\_single\_mode.m}
  \item \texttt{simulations/fem/FEM\_modes\_norm.mat}
  \item \texttt{simulations/rom/corrected\_fem\_mode\_rom.m}
\end{itemize}

\section{From Beam-Circuit System to a Two-DOF ROM}

\subsection{Beam Moment and Governing Equation}

The bending moment is modeled as the elastic moment plus two localized
piezoelectric moment contributions:
\begin{equation}
  M =
  -YI\frac{\partial^2 w_{\mathrm{rel}}}{\partial x^2}
  +\vartheta v_1(t)
    \left[\Heav(x)-\Heav\!\left(x-\frac{l}{2}\right)\right]
  +\vartheta v_2(t)
    \left[\Heav\!\left(x-\frac{l}{2}\right)-\Heav(x-l)\right].
\end{equation}
Here $v_1$ is the harvester-path voltage and $v_2$ is the modulator-path
voltage.  The bending stiffness and piezoelectric coupling coefficient used in
the archived notes are
\begin{align}
  YI &=
  \frac{2b}{3}\left[
    Y_s\frac{h_s^3}{8}
    +c_{11}^{E}\left(\left(h_p+\frac{h_s}{2}\right)^3-\frac{h_s^3}{8}\right)
  \right],\\
  \vartheta_{\mathrm{series}}
  &= \frac{e_{31}b}{2}(h_p+h_s).
\end{align}

Neglecting beam damping in the main derivation gives
\begin{equation}
  -\frac{\partial^2 M}{\partial x^2}
  +m\frac{\partial^2 w_{\mathrm{rel}}}{\partial t^2}
  =
  -m\ddot w_b(t).
\end{equation}

\subsection{Single-Mode Projection}

Projecting onto the first mass-normalized mode gives
\begin{equation}
  \ddot \eta_1+\omega_1^2\eta_1-\chi_{11}v_1-\chi_{12}v_2=f_1(t),
  \label{eq:dim_mech_voltage}
\end{equation}
where
\begin{align}
  \omega_1^2
  &=\int_0^l YI\,\phi_1^{(4)}(x)\phi_1(x)\,\dd x,\\
  \chi_{11}
  &=\vartheta\int_0^{l/2}\phi_1''(x)\,\dd x
    =\vartheta\phi_1'\!\left(\frac{l}{2}\right),\\
  \chi_{12}
  &=\vartheta\int_{l/2}^{l}\phi_1''(x)\,\dd x
    =\vartheta\left[\phi_1'(l)-\phi_1'\!\left(\frac{l}{2}\right)\right],\\
  f_1(t)
  &=-m\ddot w_b(t)\int_0^l \phi_1(x)\,\dd x.
\end{align}

The analytic plain-cantilever mode used in the initial notes is
\begin{equation}
  \phi_r(x)=\frac{1}{\sqrt{ml}}\left[
  \cos\frac{\lambda_r x}{l}-\cosh\frac{\lambda_r x}{l}
  +\frac{\sin\lambda_r-\sinh\lambda_r}{\cos\lambda_r+\cosh\lambda_r}
  \left(\sin\frac{\lambda_r x}{l}-\sinh\frac{\lambda_r x}{l}\right)
  \right],
\end{equation}
with
\begin{equation}
  1+\cos\lambda_r\cosh\lambda_r=0.
\end{equation}
In the corrected model, $\phi_1$, $\omega_1$, $\chi_{11}$, $\chi_{12}$, and
$\int_0^l\phi_1\,\dd x$ should be taken from the FEM-informed mode data rather
than from the plain beam formula.

\subsection{Harvester Circuit}

Let $\lambda_1$ be the flux linkage in the harvester path so that
$v_1=\dot\lambda_1$.  The harvester circuit equation is
\begin{equation}
  C_p\ddot\lambda_1
  +\frac{1}{R_1}\dot\lambda_1
  +\frac{1}{L}\lambda_1
  +\frac{1}{L_c}\lambda_1^3
  +\chi_{11}\dot\eta_1
  =0.
  \label{eq:dim_circuit}
\end{equation}

Using $v_1=\dot\lambda_1$, the mechanical equation becomes
\begin{equation}
  \ddot \eta_1+\omega_1^2\eta_1-\chi_{11}\dot\lambda_1
  =
  \chi_{12}v_2(t)+f_1(t).
  \label{eq:dim_mech_flux}
\end{equation}

Equations~\eqref{eq:dim_circuit} and \eqref{eq:dim_mech_flux} form the
two-degree-of-freedom reduced system.  The two solved DOFs are the structural
modal coordinate $\eta_1$ and the harvester flux linkage $\lambda_1$.  The
modulator signal $v_2(t)$ is prescribed.

\section{Normalization}

The archived notes use
\begin{equation}
  \lambda_0=\sqrt{L_cC_p\omega_1^2},\qquad
  \eta_0=\lambda_0\sqrt{C_p},\qquad
  t_0=\frac{1}{\omega_1},\qquad
  T=\frac{t}{t_0}.
\end{equation}
Define dimensionless variables
\begin{equation}
  \Lambda_1=\frac{\lambda_1}{\lambda_0},
  \qquad
  H_1=\frac{\eta_1}{\eta_0}.
\end{equation}
The dimensionless parameters are
\begin{equation}
  \zeta=\frac{1}{C_pR_1\omega_1},\qquad
  \Omega_t=\frac{\omega_t}{\omega_1},\qquad
  \omega_t^2=\frac{1}{C_pL},\qquad
  \Theta=\frac{\chi_{11}}{\omega_1\sqrt{C_p}}.
\end{equation}

For harmonic base excitation and prescribed modulation voltage,
\begin{equation}
  w_b(t)=w_0 e^{\jj\omega_b t}+\cc,
  \qquad
  v_2(t)=v_0 e^{\jj\omega_v t}+\cc,
\end{equation}
the normalized forcing amplitudes are
\begin{align}
  F_b
  &=
  \frac{m\omega_b^2w_0}{\omega_1^2\eta_0}
  \int_0^l\phi_1(x)\,\dd x,\\
  F_v
  &=
  \frac{\Theta\chi_{12}v_0}{\chi_{11}\omega_1\lambda_0}.
\end{align}

The main normalized two-DOF ROM is therefore
\begin{align}
  H_1''+H_1-\Theta\Lambda_1'
  &=
  F_b e^{\jj\Omega_bT}+F_v e^{\jj\Omega_vT}+\cc,
  \label{eq:nondim_mech}\\
  \Lambda_1''
  +\zeta\Lambda_1'
  +\Omega_t^2\Lambda_1
  +\Lambda_1^3
  +\Theta H_1'
  &=0,
  \label{eq:nondim_circuit}
\end{align}
where primes denote derivatives with respect to $T$ and
$\Omega_b=\omega_b/\omega_1$, $\Omega_v=\omega_v/\omega_1$.

\paragraph{Optional structural damping.}
For damping analysis, the structural equation is extended to
\begin{equation}
  H_1''+2\xi H_1'+H_1-\Theta\Lambda_1'
  =
  F e^{\jj\Omega_fT}+\cc.
  \label{eq:damped_mech}
\end{equation}

\section{Slow-Flow Analysis for the Two-DOF System}

\subsection{Slowly Varying Effective Forcing}

When the base and modulator frequencies are close, the forcing can be viewed as
a carrier at $\Omega_f\approx \Omega_b$ with a slowly varying envelope:
\begin{equation}
  F_b e^{\jj\Omega_bT}+F_v e^{\jj\Omega_vT}
  =
  \left(F_b+F_v e^{\jj(\Omega_v-\Omega_b)T}\right)e^{\jj\Omega_bT}.
\end{equation}
For the frozen-envelope analysis, write this as
\begin{equation}
  F e^{\jj\Omega_fT}+\cc,
\end{equation}
where $F$ is treated as slowly varying relative to the carrier oscillation.

\subsection{Complex Amplitude Variables}

Define the complex amplitude variables
\begin{equation}
  \psi_H = H_1' + \jj\Omega_f H_1,\qquad
  \psi_\Lambda = \Lambda_1' + \jj\Omega_f \Lambda_1.
\end{equation}
Then
\begin{equation}
  H_1'=\frac{\psi_H+\overline{\psi_H}}{2},\qquad
  H_1=\frac{\psi_H-\overline{\psi_H}}{2\jj\Omega_f},
\end{equation}
and similarly for $\Lambda_1$.  Assume
\begin{equation}
  \psi_H=\varphi_H e^{\jj\Omega_fT},\qquad
  \psi_\Lambda=\varphi_\Lambda e^{\jj\Omega_fT},
\end{equation}
where $\varphi_H$ and $\varphi_\Lambda$ vary slowly.

Keeping only the resonant $e^{\jj\Omega_fT}$ terms gives the slow-flow system
\begin{align}
  \dot\varphi_H
  +\frac{1-\Omega_f^2}{2\jj\Omega_f}\varphi_H
  -\frac{\Theta}{2}\varphi_\Lambda
  &=F,
  \label{eq:slow_h}\\
  \dot\varphi_\Lambda
  +\frac{\zeta}{2}\varphi_\Lambda
  +\frac{\Omega_t^2-\Omega_f^2}{2\jj\Omega_f}\varphi_\Lambda
  -\frac{3}{8\Omega_f^3}\jj |\varphi_\Lambda|^2\varphi_\Lambda
  +\frac{\Theta}{2}\varphi_H
  &=0.
  \label{eq:slow_l}
\end{align}

With structural damping $\xi$, equation~\eqref{eq:slow_h} becomes
\begin{equation}
  \dot\varphi_H
  +\left(\xi+\frac{1-\Omega_f^2}{2\jj\Omega_f}\right)\varphi_H
  -\frac{\Theta}{2}\varphi_\Lambda
  =F.
  \label{eq:slow_h_damped}
\end{equation}

\section{Steady-State Cubic and Hysteresis}

\subsection{No Structural Damping}

At steady state, eliminate $\varphi_H$ from
\eqref{eq:slow_h}--\eqref{eq:slow_l}.  The result can be written
\begin{equation}
  A\varphi_\Lambda+C|\varphi_\Lambda|^2\varphi_\Lambda=B,
  \label{eq:complex_cubic}
\end{equation}
with
\begin{align}
  A
  &=
  \frac{\zeta}{2}\frac{1-\Omega_f^2}{2\jj\Omega_f}
  -\frac{(\Omega_t^2-\Omega_f^2)(1-\Omega_f^2)}{4\Omega_f^2}
  +\frac{\Theta^2}{4},\\
  B&=-\frac{\Theta}{2}F,\\
  C&=-\frac{3(1-\Omega_f^2)}{16\Omega_f^4}.
\end{align}

Taking the squared magnitude gives the real cubic relation
\begin{equation}
  C^2\rho^3
  +2\operatorname{Re}(A)C\rho^2
  +|A|^2\rho
  =
  |B|^2,
  \qquad
  \rho=|\varphi_\Lambda|^2.
  \label{eq:real_cubic}
\end{equation}
This is the algebraic equation used to plot the hysteresis curve and determine
stable/unstable branches.

\subsection{With Structural Damping}

With structural damping, define
\begin{equation}
  D=\xi+\frac{1-\Omega_f^2}{2\jj\Omega_f}.
\end{equation}
The same cubic form~\eqref{eq:complex_cubic} holds, but with
\begin{align}
  A
  &=
  \frac{\zeta}{2}D
  +\frac{\Omega_t^2-\Omega_f^2}{2\jj\Omega_f}D
  +\frac{\Theta^2}{4},\\
  B&=-\frac{\Theta}{2}F,\\
  C&=-\frac{3}{8\Omega_f^3}\jj D.
\end{align}

\section{Linear Stability of a Slow-Flow Branch}

Let a steady solution be
$(\varphi_{H,0},\varphi_{\Lambda,0})$.  Perturb it as
\begin{equation}
  \varphi_H=\varphi_{H,0}+\delta_H,\qquad
  \varphi_\Lambda=\varphi_{\Lambda,0}+\delta_\Lambda.
\end{equation}
The perturbation dynamics can be written in the standard complex-conjugate
linear form
\begin{equation}
  \frac{\dd}{\dd T}
  \begin{bmatrix}
    \delta_H\\
    \delta_\Lambda
  \end{bmatrix}
  =
  \bm A
  \begin{bmatrix}
    \delta_H\\
    \delta_\Lambda
  \end{bmatrix}
  +
  \bm B
  \begin{bmatrix}
    \overline{\delta_H}\\
    \overline{\delta_\Lambda}
  \end{bmatrix}.
\end{equation}
For the undamped structural case,
\begin{align}
  \bm A
  &=
  -\begin{bmatrix}
  \dfrac{1-\Omega_f^2}{2\jj\Omega_f}
  &
  -\dfrac{\Theta}{2}\\[6pt]
  \dfrac{\Theta}{2}
  &
  \dfrac{\zeta}{2}
  +\dfrac{\Omega_t^2-\Omega_f^2}{2\jj\Omega_f}
  -\dfrac{3\jj}{4\Omega_f^3}
  |\varphi_{\Lambda,0}|^2
  \end{bmatrix},\\
  \bm B
  &=
  -\begin{bmatrix}
  0&0\\
  0&-\dfrac{3\jj}{8\Omega_f^3}\varphi_{\Lambda,0}^2
  \end{bmatrix}.
\end{align}
The branch is stable if all eigenvalues of
\begin{equation}
  \begin{bmatrix}
    \bm A & \bm B\\
    \overline{\bm B} & \overline{\bm A}
  \end{bmatrix}
\end{equation}
have non-positive real part.  This is the stability check implemented in the
ROM scripts when marking unstable portions of the hysteresis curve.

\section{Power and Energy Accounting}

The harvester output power is
\begin{equation}
  P_1(t)=\frac{v_1^2(t)}{R_1}
  =
  \frac{1}{R_1}\dot\lambda_1^2.
\end{equation}
In normalized variables, it is proportional to
\begin{equation}
  P_{1,\mathrm{norm}}=\zeta\left(\Lambda_1'\right)^2.
\end{equation}

The prescribed modulator path does work on the beam through
\begin{equation}
  P_2(t)=v_2(t)i_2(t),
\end{equation}
where
\begin{equation}
  i_2(t)=C_p\dot v_2(t)+\chi_{12}\dot\eta_1(t).
\end{equation}
The capacitive term is reversible charging energy.  The mechanically effective
term is therefore
\begin{equation}
  P_{2,\mathrm{eff}}(t)=\chi_{12}v_2(t)\dot\eta_1(t).
\end{equation}

The January 31 slides and \texttt{simulations/rom/corrected\_fem\_mode\_rom.m}
compare cumulative normalized harvested energy against modulation input.  This
is the right diagnostic for judging whether the nonlinear branch transition is
useful, rather than merely producing a larger response amplitude.

\section{Relation to Simulation Results}

\begin{itemize}
  \item \textbf{Initial single-mode ROM.}  Hysteresis and branch selection under
  harmonic base excitation and modulation were checked in the older analytic
  mode scripts in \path{simulations/legacy_reference/}.
  \item \textbf{Continuous beam FEM.}  The full beam model was used to test
  whether the ROM-predicted branch transition appears beyond the single-mode
  approximation.  See \path{simulations/fem/continuous_harvester.m}.
  \item \textbf{Corrected mode assumption.}  The FEM-informed mode shape and
  frequency replace the plain cantilever mode.  See
  \path{simulations/fem/eigenanalysis_corrected_single_mode.m}.
  \item \textbf{Corrected ROM.}  The main two-DOF model loads
  \path{FEM_modes_norm.mat} and should be compared against FEM.  See
  \path{simulations/rom/corrected_fem_mode_rom.m}.
  \item \textbf{Power flow.}  Harvested energy versus modulation input is
  discussed in the January 31 slides and implemented in the corrected ROM
  scripts under \path{simulations/rom/}.
  \item \textbf{Second-harmonic robustness.}  Nonlinear and linear comparisons
  with a secondary base-excitation harmonic are archived under
  \path{simulations/robustness/second_harmonics/}.
\end{itemize}

\section{Open Points for Collaborators}

\begin{itemize}
  \item The sign convention for $\chi_{11}$, $\chi_{12}$, and the FEM mode
  orientation should be kept consistent.  The corrected MATLAB scripts use the
  negative of the FEM mode quantities loaded from \texttt{FEM\_modes\_norm.mat}.
  \item The current two-DOF ROM treats the modulator voltage as prescribed.  If
  the modulator electronics are later modeled as a real circuit, it will add
  additional electrical states.
  \item The multi-mode Word note and \texttt{Old\_sources/Multi\_DOF/} are
  useful future directions, but the present mainline result is still the
  corrected single-mode ROM.
  \item Brownian or rough stochastic base motion needs care because the FEM
  boundary formulation uses prescribed acceleration.  This was already flagged
  in the February 14 meeting.
\end{itemize}

\end{document}
